\section{Discussion}
\label{sec:discussion}

\textbf{A recall number is an operating point.} Recall and false positives move together with any
severity threshold, and a critic can be tuned to any recall by blocking more. On general-purpose code
the same auditor sits at different points under different decision rules, repetition buys recall at
a roughly constant price in false positives, and pooling families buys coverage at a higher one. On
scientific code the operating point also moved with what the task named as the deliverable. A recall
figure quoted without its false-positive rate, its rulebook, its decision rule and its task framing
does not carry to another setting.

\textbf{What the science studies add.} They give the science profile's design rule (the model names
evidence, code verifies it) an empirical reading and a boundary. On results, number$\to$source
verified every inconsistency between a report, its results file and its log, and could not see a
result that was wrong consistently; the model, by recomputing, flagged 23 of 25 such fabrications. On
data, a validator written from the card caught what a rule can state and missed what needs reading
across columns and units, which the model caught in 18 of the 28 such items the validator missed.
Where both tiers ran, the deterministic one flagged no clean item and the model carried the errors;
on data the union was better than either part, and on results the model's flags already contained
the profile's, whose findings its prompt carried. Neither study shows that the model tier scales: the
artefacts were small, the model's fabrication findings are hand calculations (several wrong), and it
missed every fault that needs comparison across many rows. A model reading is a useful second tier
where the artefact fits in the reader's working memory, not a substitute for executing the code or
computing over the data.

\textbf{Clean items must be clean by the auditor's standard.} In all three science studies the auditor
flagged defects in our own construction: a card and file that named a column differently, a file
outside the audited scope, a report claiming a run no file performed, a provenance sentence false for
inputs passed as literals, a unit label contradicting the code's documentation, and a task that named
the whole problem when one step was delivered. Each was a correct flag on an item we had built as
clean. We caught two in pilots, one when every verdict came back as an escalation, and the rest only
by reading every flag on a clean item, which the registrations required. Without that reading, the
false-positive rate would have measured the evaluator.

\textbf{Limits.} The auditor families are not sampled alike, which bounds every cross-family union
contrast, and no family was compared at a matched false-positive rate. Primary intervals are
problem-cluster bootstraps whose coverage was simulated only for single rates. Every rating was made
by a model, one of them the agent that designed the studies; no human rated anything. The reviewer of
every report is also one of the measured families. The science studies use one benchmark or four
tables each, synthetic faults, small artefacts and one generator. The full limits, and the
experiments that would remove them, are in Appendices~\ref{sec:discussion-full} and~\ref{sec:limitations-full}.
